\documentclass{article}
\begin{document}
\section*{Independent review of cartesian preservation by native flattening}
A clean Lean 4.35.0 build ran 800 jobs and 1,692 indexed declaration axiom audits. Sixty previous project modules are unchanged. Author acceptance probes were replayed. The review independently downloaded the candidate from Library and checked its digest.
\par Weight208Review constructs a weighted arrow over $2\to5$ increasing a terminal payload from 3 to 9. A cartesian factor would force $9\le3$. Its image stays noncartesian, but still has a unique incoming lift at the specified target. Weighted flattening need not be full: the outer empty and single empty block have the same weighted image and no cross-tag arrow.
\par Gauge208Review and WeightedGauge208Review use different nonnormal grade choices. The labels for the comparisons over $2\to5$, $5\to7$ and $2\to7$ are respectively 6, 32, and 8; the identity comparison at 5 has label 24. Naive addition of the first two is not the composition law. The proved laws include the source and target choices' composition and unit cells. The weighted probe instantiates these laws and the mixed equation with a nonfull constant payload transformation.
\par WeakBoundary208Review exhibits a positive-length equipment comparison whose functor does not have strict unique incoming lifts. Its nonliteral boundary object is only isomorphic to an image. Thus the strict matching model cannot be replaced without transport by an equivalent iso-comma model.
\par The weighted cartesian theorem is relative to the source projection. The weighted comparison uses independently chosen grade cleavages and their specified payload extensions, not arbitrary independently chosen total cleavages. Three-level associativity, action on maps of spans, identification with the packed multiplication and strict contextad recovery remain open. No blocker found within this scope.
\end{document}
